\subsection{Elementary tests of the mathematical distinctions}
\label{subsec:contrastive-tests}
The following elementary examples make the comparison concrete. Each
isolates a boundary between a useful intermediate statement and a stronger
conclusion required by one of the reviewed arguments.

\begin{example}[Restricted energy control and the full spectral gap]
On $\{-1,1\}^2$ with uniform law, consider the continuous-time chain
that flips the first coordinate at rate $1$ and the second at rate
$\tau>0$. Its Dirichlet form is
\[
 \mathcal E_\tau(f)=\frac12\mathbb E\!\left[
 (f(-a,b)-f(a,b))^2+\tau(f(a,-b)-f(a,b))^2\right].
\]
For every observable $f=g(a)$, one has
$\mathcal E_\tau(f)=2\operatorname{Var}(f)$, uniformly in $\tau$.
For $f(a,b)=b$, however, the variance is $1$ and the energy is $2\tau$.
The functions $1,a,b,ab$ diagonalize the negative generator, with
eigenvalues $0,2,2\tau,2(1+\tau)$. Thus the full spectral gap is
$2\min\{1,\tau\}$, although the restricted inequality remains fixed.
At $\tau=0$ the restricted inequality even coexists with reducibility.
This locates the exact distinction between the matching argument's
upgrade to a full inequality and the common-base argument's separate
use of selected observables and return control. A restricted estimate
must be paired with the additional property the algorithm actually uses.
\end{example}

\begin{example}[The dimension loss in an entrywise operator estimate]
Suppose the scalar inequality
$\|p(A)\|\le C\max_{z\in W(A)}|p(z)|$ is available, and let
$Q(z)=[q_{ij}(z)]_{i,j=1}^m$ be a matrix polynomial. Viewing $Q[A]$
as a block matrix gives the elementary estimate
\[
 \|Q[A]\|\le m\max_{i,j}\|q_{ij}(A)\|
 \le mC\max_{z\in W(A)}\|Q(z)\|.
\]
Indeed, bound each output block by the largest block norm times
$\sum_j\|x_j\|$, then apply Cauchy--Schwarz to input and output blocks.
This argument retains a factor $m$. The complete Crouzeix claim seeks
constant $2$ independently of $m$; its ordered product tests must remove
the dimensional loss that this entrywise route leaves behind. The
distinction concerns the strength of the argument, even when its scalar
inputs are already sharp.
\end{example}

\begin{example}[Positive nodes and a continuum inequality]
For $h>0$, the polynomial
$r_h(t)=(t-h/2)^2-h^2/16$ is positive at every node $t\in h\mathbb Z$,
where $r_h(t)\ge3h^2/16$, but $r_h(h/2)=-h^2/16$.
A positive node certificate therefore needs a separate argument between
nodes. For example, on a covered compact interval, a derivative bound
$|r'|\le L$ and node lower bounds $r\ge Lh/2$ imply $r\ge0$ whenever
every point is within $h/2$ of a node. An all-real conclusion additionally
needs control beyond the covered interval. This explains the distinct
roles of rational node enclosures, interpolation, and tail estimates in
the general Mahler manuscript. The proposed advance is their compatibility
with the scalar error budget needed by the matrix comparison.
\end{example}
